% Generated draft: requires mathematical and source review.
\documentclass[11pt]{article}
\usepackage[margin=25mm]{geometry}
\usepackage{fontspec}
\setmainfont{Latin Modern Roman}
\usepackage{amsmath,amssymb,mathtools,unicode-math}
\setmathfont{Latin Modern Math}
\usepackage{enumitem,needspace,xurl,hyperref}
\setlength{\parindent}{0pt}
\setlength{\parskip}{5pt}
\setlength{\emergencystretch}{3em}
\newcommand{\N}{\mathbb{N}}
\newcommand{\Z}{\mathbb{Z}}
\newcommand{\Q}{\mathbb{Q}}
\newcommand{\R}{\mathbb{R}}
\newcommand{\C}{\mathbb{C}}
\newcommand{\F}{\mathbb{F}}
\newcommand{\A}{\mathbb{A}}
\newcommand{\PP}{\mathbb{P}}
\newcommand{\E}{\mathbb{E}}
\newcommand{\eps}{\varepsilon}
\newcommand{\dd}{\,\mathrm{d}}
\begin{document}
\section*{UnsolvedMath ID: 30002123}
\textbf{Source problem code:} catalogue entry
\subsection{Definitions and mathematical statement}
A finite geometric simplicial complex \(C\) in \(\R^n\) is a finite collection of Euclidean simplices, closed under taking faces, such that the intersection of any two simplices is empty or a common face. Its underlying space \(|C|\) is the union of its simplices. Under the geometric interpretation of ``convex,'' \(|C|\) is convex: for every \(x,y\in |C|\) and every \(t\in[0,1]\), one has \((1-t)x+ty\in |C|\).

The barycentric subdivision \(\operatorname{sd} C\) has a vertex \(b_\sigma\), the barycenter of \(\sigma\), for each nonempty face \(\sigma\) of \(C\). Its nonempty simplices correspond to strictly increasing chains
\[
\sigma_0\subsetneq\sigma_1\subsetneq\cdots\subsetneq\sigma_k
\]
of nonempty faces of \(C\), with associated vertices \(b_{\sigma_0},\ldots,b_{\sigma_k}\).

A facet is an inclusion-maximal face. A finite simplicial complex is pure of dimension \(d\) if every facet has dimension \(d\). For \(d\geq1\), a finite pure \(d\)-dimensional complex \(K\) is shellable if its facets admit an ordering \(F_1,\ldots,F_m\) such that, for every \(j\in\{2,\ldots,m\}\), the subcomplex
\[
\langle F_j\rangle\cap\bigcup_{i<j}\langle F_i\rangle
\]
is nonempty and pure of dimension \(d-1\). Here \(\langle F\rangle\) denotes the simplex \(F\) together with all its faces. A nonempty zero-dimensional complex is shellable by convention.

\textbf{Question, under the finite geometric interpretation.} For every integer \(n\geq0\) and every nonempty finite geometric simplicial complex \(C\) in \(\R^n\) whose underlying space is convex, is \(\operatorname{sd} C\) shellable? The dimension is the intrinsic dimension of \(|C|\); full dimension in \(\R^n\) is not required. Under this interpretation, \(C\) is pure. The question concerns the first barycentric subdivision of the entire complex, including its boundary faces.

\textbf{Unresolved source qualification.} The supplied statement says only ``convex simplicial complex.'' It does not explicitly specify finiteness, nonemptiness, a geometric realization with convex underlying space, or a convention for shellability. The formulation above makes a standard finite geometric reading explicit; confirmation that these conventions match the report's intended scope requires human review.
\subsection{Short English statement}
Does taking the first barycentric subdivision make every finite simplicial triangulation of a nonempty convex set shellable? This expresses the supplied question under an explicit geometric interpretation of ``convex simplicial complex''; the source's intended conventions still require confirmation.
\subsection{Sources}
\begin{itemize}
\item Local registry: \texttt{data/unsolvedmath/problems.json}, permanent UnsolvedMath ID \(30002123\), source problem code \texttt{catalogue entry}, title ``Shellability of Barycentric Subdivisions of Convex Complexes,'' in the Oberwolfach Reports Open Problems collection. The supplied record lists proposed year 2012 and supplies no named proposer.
\item Supplied report citation: \emph{Discrete Differential Geometry (2013)}, Oberwolfach Reports, DOI 10.4171/owr/2012/34, pp. 2077--2137. \url{https://doi.org/10.4171/owr/2012/34}. The citation's title year 2013 and the record's proposed year 2012 are retained as supplied; their relationship has not been independently checked.
\item Supplied literature source, labeled ``Adiprasito Benedetti'': \url{https://arxiv.org/abs/1709.07930}.
\item Supplied literature source, labeled ``OWR report'': \url{https://www.ihes.fr/~fguenth/files/OWR_2012_34.pdf}.
\item Supplied literature assessment dated \texttt{2026-08-21}: the record labels the problem open and reports that general shellability of the first barycentric subdivision remains unresolved, while attributing shellability after a bounded further subdivision and related collapsibility results to Adiprasito and Benedetti. This is supplied dated curation derived from the OWR collection, not an independently verified literature assessment; its status and descriptions of cited results require independent verification.
\item Supplied statement-verification metadata: \texttt{corrected\_verified}. The record reports that mathematical content was checked against the official Oberwolfach report and that notation, context, and question form were normalized without changing the source-level claim. This is a supplied verification claim, not verification performed for this draft.
\end{itemize}
\end{document}
